\documentclass[12pt,a4paper]{report}
\usepackage[utf8]{inputenc}
\usepackage[T1]{fontenc}
\usepackage{amsmath, amsthm, amssymb,latexsym,amscd,amsfonts,enumerate}
\usepackage[top=2.5cm, bottom=2.0cm, left=1.5cm, right=1.5cm]{geometry}
\usepackage{fancyhdr, graphicx, wrapfig}
\pagestyle{fancy}

% 1. Setup for 'fancy' page style (standard pages)
\fancyhf{} 
\rfoot{Page \thepage} 
\renewcommand{\headrulewidth}{0pt} 
\renewcommand{\footrulewidth}{1.2pt} 

% 2. Redefine 'plain' style (ToC, Chapter starts)
\fancypagestyle{plain}{%
  \fancyhf{}
  \rfoot{Page \thepage}
  \renewcommand{\headrulewidth}{0pt}
  \renewcommand{\footrulewidth}{1.2pt}
}

% 3. Redefine 'empty' style (Cover page)
\fancypagestyle{empty}{%
  \fancyhf{}
  \rfoot{Page \thepage}
  \renewcommand{\headrulewidth}{0pt}
  \renewcommand{\footrulewidth}{1.2pt}
}

\newcommand{\tab}{\hspace{0.5 cm}}
\usepackage{xcolor}
\usepackage{tikz}
\usepackage{scrextend}
\usepackage{algorithm}
\usepackage{algpseudocode}
\pagestyle{plain}
\pagestyle{fancy}
\lhead{\it Major Project Proposal}
\rhead{\it Computer Vision: Product Tracking}
\lfoot{\it Group: Do Duy Loi, Trinh Chan Duy \& Dang Vo Hong Phuc}
\rfoot{Page \thepage}

\fancypagestyle{plain}{%
  \fancyhf{} 
  \lhead{\it Major Project Proposal}
  \rhead{\it Computer Vision: Product Tracking}
  \lfoot{\it Group: Do Duy Loi, Trinh Chan Duy \& Dang Vo Hong Phuc}
  \rfoot{Page \thepage}
  \renewcommand{\headrulewidth}{1.2pt}
  \renewcommand{\footrulewidth}{1.2pt}
}

\renewcommand{\headrulewidth}{1.2pt}
\renewcommand{\footrulewidth}{1.2pt}
\usepackage{booktabs}
\usepackage{array}
\usepackage{caption}
\usepackage{forest}
\usepackage{mdframed}
\usepackage{float}
\usepackage[english]{babel}
\usepackage{tocloft}
\usepackage{helvet}
\usepackage{tabularx}
\usepackage[unicode, colorlinks=true,linkcolor=black,urlcolor=blue]{hyperref} 

% Redefine Table of Contents behavior
\makeatletter
\renewcommand\tableofcontents{%
    \begin{center}
      {\Large\bfseries \contentsname} 
      \par\nobreak
      \vspace{10pt} 
    \end{center}
    \@starttoc{toc}
}
\makeatother
\renewcommand{\familydefault}{\sfdefault}

\begin{document}
\fontsize{13pt}{18pt}\selectfont
\setlength{\baselineskip}{18truept}

\begin{titlepage}
\begin{tikzpicture}[remember picture,overlay,inner sep=0,outer sep=0]
    \draw[blue!70!black,line width=4pt] ([xshift=-1.5cm,yshift=-2cm]current page.north east) coordinate (A)--([xshift=1.5cm,yshift=-2cm]current page.north west) coordinate(B)--([xshift=1.5cm,yshift=2cm]current page.south west) coordinate (C)--([xshift=-1.5cm,yshift=2cm]current page.south east) coordinate(D)--cycle;
    \draw ([yshift=0.5cm,xshift=-0.5cm]A)-- ([yshift=0.5cm,xshift=0.5cm]B)--
    ([yshift=-0.5cm,xshift=0.5cm]B) --([yshift=-0.5cm,xshift=-0.5cm]B)--([yshift=0.5cm,xshift=-0.5cm]C)--([yshift=0.5cm,xshift=0.5cm]C)--([yshift=-0.5cm,xshift=0.5cm]C)-- ([yshift=-0.5cm,xshift=-0.5cm]D)--([yshift=-0.5cm,xshift=-0.5cm]D)--([yshift=0.5cm,xshift=0.5cm]D)--([yshift=-0.5cm,xshift=0.5cm]A)--([yshift=-0.5cm,xshift=-0.5cm]A)--([yshift=0.5cm,xshift=-0.5cm]A);
    \draw ([yshift=-0.3cm,xshift=0.3cm]A)-- ([yshift=-0.3cm,xshift=-0.3cm]B)--
    ([yshift=0.3cm,xshift=-0.3cm]B) --([yshift=0.3cm,xshift=0.3cm]B)--([yshift=-0.3cm,xshift=0.3cm]C)--([yshift=-0.3cm,xshift=-0.3cm]C)--([yshift=0.3cm,xshift=-0.3cm]C)-- ([yshift=0.3cm,xshift=0.3cm]D)--([yshift=-0.3cm,xshift=0.3cm]D)--([yshift=-0.3cm,xshift=-0.3cm]D)--([yshift=0.3cm,xshift=-0.3cm]A)--([yshift=0.3cm,xshift=0.3cm]A)--([yshift=-0.3cm,xshift=0.3cm]A);
\end{tikzpicture}
\begin{center}
{\large\bf UNIVERSITY OF SCIENCE}\\
{---------------------o0o--------------------}\\
\vskip 1cm
{\bf ACADEMIC COURSE REPORT\\ MAJOR PROJECT PROPOSAL}\\[0.5 cm]
\rule{\linewidth}{0.5 mm} \\[0.2 cm]
{\Large\bf \textbf{PRODUCT DETECTION AND TRACKING\\ ON MANUFACTURING CONVEYOR BELTS}}\\
\rule{\linewidth}{0.5 mm} \\
\vskip 1.5cm
\includegraphics[scale=0.25]{format/logohcmus.jpg}\\[1cm] 
\begin{minipage}{0.45\textwidth}
\begin{flushleft} \large
\emph{{\bf Group Members:}}\\
{\bf Do Duy Loi}\\ 
{\bf Trinh Chan Duy}\\
{\bf Dang Vo Hong Phuc}
\end{flushleft}
\end{minipage}
~
\begin{minipage}{0.4\textwidth}
\begin{flushright} \large
\emph{{\bf Student ID:}} \\
{\bf 23120293}\\
{\bf 23120419}\\
{\bf 23120155}
\end{flushright}
\end{minipage}\\
\vfill
{\bf Ho Chi Minh City, April 2026}
\end{center}
\end{titlepage}

\tableofcontents
\newpage
\pagenumbering{arabic}

\chapter{Project Team Information}
\renewcommand{\arraystretch}{1.2}

\section{Members}
\begin{center}
\begin{tabularx}{\textwidth}{l c X}
\toprule
\textbf{Full Name} & \textbf{Student ID}  \\ 
\midrule
Do Duy Loi & 23120293  \\
Trinh Chan Duy   & 23120419  \\ 
Dang Vo Hong Phuc   & 23120155  \\
\bottomrule
\end{tabularx}
\end{center}

\section{Project Title}
\textbf{Computer Vision System for Automated Product Detection, Tracking, and Monitoring on Manufacturing Conveyor Belts.}

\chapter{Project Overview}

\section{Research Motivation}
In the era of Industry 4.0, automating production processes is critical for maintaining corporate competitiveness. Monitoring product flow on conveyor belts is a vital stage in this automation. Currently, many factories still rely on manual labor for counting and inspection, leading to fatigue-induced errors, slow processing speeds, and high operational costs.

There is a practical demand for a system capable of 24/7 operation with high precision to:
\begin{itemize}
    \item Automate the counting of output products.
    \item Monitor flow to detect early-stage blockages or misaligned items.
    \item Collect operational data to optimize assembly line efficiency.
\end{itemize}

\section{Problem Statement}
\begin{itemize}
    \item \textbf{Input:} Real-time video stream from industrial cameras mounted above the conveyor belt, observing products such as bottles, cartons, or mechanical parts.
    \item \textbf{Output:} 
    \begin{itemize}
        \item Bounding boxes identifying the location of each product.
        \item Unique IDs assigned to each product via multi-object tracking.
        \item Trajectories and status indicators (Normal/Defective).
        \item Statistical reports on total product throughput within specific timeframes.
    \end{itemize}
\end{itemize}
\newpage
\section{Scientific and Practical Significance}
\begin{itemize}
    \item \textbf{Scientific Significance:} This project evaluates the performance of state-of-the-art Deep Learning architectures (YOLO) combined with data filtering algorithms (Kalman Filter, Re-ID) in specific industrial environments. It explores model optimization to achieve a balance between accuracy and real-time processing speed.
    \item \textbf{Practical Significance:} It minimizes human error, increases assembly line productivity by 20-30\%, provides accurate digital data for Manufacturing Execution Systems (MES/ERP), and serves as a foundation for automated Quality Assurance (QA/QC).
\end{itemize}

\section{Challenges and Scope}
\subsection{Technical Challenges}
\begin{itemize}
    \item \textbf{Conveyor Speed:} High-speed movement can cause motion blur.
    \item \textbf{Environmental Variance:} Unstable factory lighting and industrial dust accumulation on lenses.
    \item \textbf{Product Similarity:} Identical appearance of products makes unique ID assignment difficult during occlusion.
    \item \textbf{Product Density:} Crowded or overlapping items can confuse detection models.
\end{itemize}

\subsection{Project Scope}
\begin{itemize}
    \item \textbf{Datasets:} Standard benchmarks including MVTec D2S (instance segmentation), MVTec AD (anomaly detection), and custom data collected via Roboflow.
    \item \textbf{Methodology:} Focus on YOLOv8/v11 for detection and ByteTrack or Deep SORT for tracking, deployed on GPU-accelerated hardware.
\end{itemize}

\chapter{Related Research}

\section{Practical Application Studies}
Rather than focusing solely on theoretical algorithms, this project examines two real-world application systems with contexts highly relevant to industrial requirements:

\subsection{Product Quality Inspection and Counting System (MDPI, 2024)}
The research \textbf{"Computer-Vision-Based Product Quality Inspection and Novel Counting System"} presents a computer vision system deployed on an actual production line. 
\begin{itemize}
    \item \textbf{Application:} The system performs dual tasks: counting products and detecting surface defects by combining \textbf{YOLOv8 Pose} and \textbf{difference image} techniques.
    \item \textbf{Results:} Achieved \textbf{99.2\% accuracy} for products over 1kg by integrating visual data with weight sensors, operating stably at a camera frame rate of 120 FPS.
    \item \textbf{Reference:} \textit{Applied Sciences, MDPI (2024).}
\end{itemize}

\subsection{Zone-Based Product Counting System (ResearchGate, 2026)}
The study \textbf{"Real-Time Object Counting on Industrial Conveyor Belts Using YOLOv11 and Zone-Based Algorithm"} proposes a solution for conveyor belts with high object density.
\begin{itemize}
    \item \textbf{Application:} Utilizes the \textbf{YOLOv11} model combined with a \textbf{Three-zone counting} algorithm. Instead of traditional line-crossing logic, the system tracks object transitions through three defined spatial zones.
    \item \textbf{Results:} Effectively resolved the "double-counting" phenomenon and achieved \textbf{94.2\% accuracy} even under 50-70\% occlusion.
    \item \textbf{Reference:} \textit{ResearchGate (2026).}
\end{itemize}

\subsection{Comparative Analysis of Applied Research}
\begin{center}
\small
\renewcommand{\arraystretch}{1.5}
\begin{tabularx}{\textwidth}{|p{3.5cm}|X|X|}
\hline
\textbf{Research Work} & \textbf{Pros} & \textbf{Cons} \\ \hline
\textbf{MDPI System (2024)} & Extremely high accuracy (99.2\%); Multi-tasking (counting + defects); Stability through weight sensor fusion. & Requires complex hardware (120 FPS cameras, load cells); Sensitive to sudden lighting changes due to difference imaging. \\ \hline
\textbf{ResearchGate (2026)} & Handles high occlusion (50-70\%); Eliminates double-counting; YOLOv11 optimized for inference speed. & Requires precise spatial calibration; Performance depends on the stability of tracking within the central zone. \\ \hline
\end{tabularx}
\end{center}

\chapter{Proposed Approach}

\section{Identified Problems Post-Survey}
Based on the practical research surveyed above, the team has identified two core problems to resolve for system optimization:
\begin{enumerate}
    \item \textbf{Double-Counting and Omissions:} Traditional line-crossing algorithms are prone to errors when objects move inconsistently or when "flickering" occurs between frames due to environmental lighting.
    \item \textbf{System Latency and ID Switches:} Running complex Deep Learning models on high-speed lines causes inference latency, leading to deviations in the Kalman filter's predicted positions and subsequent ID swaps.
\end{enumerate}

\section{Proposed Solutions and Improvements}
To overcome the aforementioned drawbacks, the following specific improvements are proposed:
\begin{itemize}
    \item \textbf{Tracking Optimization (Inspired by YOLOv11):} Implement a \textbf{Virtual Counting Zones} logic. Instead of a single tripwire, the system will establish \textit{Entry}, \textit{Tracking}, and \textit{Exit} zones. A product is only logged into the database once it has successfully navigated all three zones in the correct temporal sequence.
    \item \textbf{Speed Optimization (Inspired by MDPI):} Utilize the \textbf{TensorRT} toolkit to optimize YOLO networks, allowing the system to maintain processing speeds above 60 FPS on Edge AI devices (e.g., Jetson Nano), ensuring trajectory continuity.
    \item \textbf{Multi-Task Defect Inspection:} Leverage data from the tracking phase to extract product images at the optimal viewpoint, then apply a lightweight classifier to detect shape abnormalities (dents, deformities) or label defects.
\end{itemize}

\chapter*{References}
\addcontentsline{toc}{chapter}{References}
\begin{enumerate}
    \item Jocher, G., et al. (2023). \textit{YOLOv8 by Ultralytics}. \url{https://github.com/ultralytics/ultralytics}.
    \item Zhang, Y., et al. (2022). \textit{ByteTrack: Multi-Object Tracking by Associating Every Detection Box}. In ECCV.
    \item Wojke, N., et al. (2017). \textit{Simple Online and Realtime Tracking with a Deep Association Metric}. In ICIP.
    \item MVTec Software GmbH. \textit{D2S Dataset for Instance Segmentation}. \url{https://www.mvtec.com/company/research/datasets/d2s/}.
    \item Bergmann, P., et al. (2019). \textit{MVTec AD Dataset for Anomaly Detection}. In CVPR.
\end{enumerate}

\end{document}
